\documentclass[11pt,a4paper]{article}

\usepackage[utf8]{inputenc}
\usepackage[main=italian,provide=*]{babel}
\usepackage{amsmath}
\usepackage{amssymb}
\usepackage{mathtools}
\usepackage{geometry}
\usepackage{booktabs}
\usepackage{hyperref}
\usepackage{seqsplit}

\geometry{margin=2.7cm}

\title{Quantum simulation del trimero a catena aperta: derivazione teorica completa\\
\large Decomposizione di Suzuki--Trotter a tre livelli per l'Hamiltoniana con termine DM}
\author{}
\date{}

\newcommand{\ket}[1]{\lvert #1 \rangle}
\newcommand{\bra}[1]{\langle #1 \rvert}

\begin{document}

\maketitle
\tableofcontents

\bigskip
\noindent Questo documento raccoglie in forma sistematica le derivazioni relative all'estensione del task di \emph{quantum simulation} dal trimero anello alla seconda topologia $N=3$: la catena aperta. Segue lo stesso schema di \texttt{\seqsplit{quantum\_simulation\_trimero\_anello\_teoria.tex}}, con due differenze strutturali che vanno dimostrate, non solo enunciate: (i) la catena ha \textbf{due} legami di scambio, non tre, e un solo commutatore da valutare al livello 1 (Sezione~\ref{sec:chirality}); (ii) l'ordine di applicazione dei due legami è \textbf{alternabile} fra passi consecutivi, una possibilità che l'anello (tre legami, non due) non ha in forma altrettanto pulita (Sezione~\ref{sec:alternanza}). Ogni identità algebrica non banale è verificata numericamente (confronto diretto di matrici $8\times8$), non solo asserita.

\section{Hamiltoniana e convenzione}
\label{sec:hamiltoniana}

Convenzione di Pauli diretta, identica nello spirito a quella dell'anello, per lo stesso motivo: il modulo \texttt{\seqsplit{trotter\_trimero\_catena.py}} costruisce l'Hamiltoniana riusando direttamente \texttt{\seqsplit{trimer\_chain\_exact.py}} come sorgente unica.

\begin{align}
H &= b\sum_{i=1}^3 Z_i \;+\; H_{ex} \;+\; H_{DM}, \label{eq:H-full}\\[4pt]
H_{ex} &= J\,\vec\sigma_1\!\cdot\!\vec\sigma_2 \;+\; J\,\vec\sigma_2\!\cdot\!\vec\sigma_3, \label{eq:Hex-def}\\[4pt]
H_{DM} &= D\,(X_1Z_2{-}Z_1X_2) \;-\; D\,(X_2Z_3{-}Z_2X_3). \label{eq:HDM-def}
\end{align}

Due differenze immediate rispetto all'anello (Eq.~1--3 di \texttt{\seqsplit{quantum\_simulation\_trimero\_anello\_teoria.tex}}): \textbf{un solo} accoppiamento di scambio $J$ (non $J,J'$: la catena non ha un terzo legame $(3,1)$ su cui distinguere basa/lati), e il termine DM ha \textbf{un solo} grado di libertà, non due opzioni A/B. La simmetria di riflessione $P_{13}$ (scambio dei siti estremi attorno al centro) impone $D_{12}=-D_{23}\equiv D$: è l'unica scelta compatibile con la simmetria, non una fra alternative (\texttt{\seqsplit{analisi\_dm\_trimero\_catena.tex}}).

\paragraph{Attenzione alla normalizzazione di $D$.} A differenza dell'anello, dove $D_{ij}=rJ_{ij}$ (il DM scala con l'accoppiamento del bond, Opzione B), qui $D$ è \textbf{assoluto}: $D_{12}=+D$, $D_{23}=-D$, senza moltiplicare per $J$. È così che \texttt{\seqsplit{trimer\_chain\_exact.py}} definisce \texttt{dm\_term}, ed è la convenzione seguita in questo documento e nel codice. Usare la convenzione dell'anello qui produrrebbe un'Hamiltoniana diversa di un fattore $J$, silenziosamente.

Si raggruppano i termini come per l'anello: $H_0=b\sum_iZ_i+H_{ex}$ (campo + scambio, fisso, indipendente da $D$) e $H_{DM}$ (si annulla per $D=0$).

\section{Decomposizione esatta di un singolo legame di scambio}
\label{sec:bond-exact}

Identica, parola per parola, alla Sezione~2 di \texttt{\seqsplit{quantum\_simulation\_trimero\_anello\_teoria.tex}}: per una coppia $(i,j)$, $J\,\vec\sigma_i\cdot\vec\sigma_j=J(X_iX_j+Y_iY_j+Z_iZ_j)$, i tre termini commutano a due a due (l'argomento non dipende da quanti altri legami esistono nel sistema), quindi
\begin{equation}
\boxed{U_{ij}(t) \equiv e^{-iJt\,\vec\sigma_i\cdot\vec\sigma_j} = R_{XX}(2Jt)\,R_{YY}(2Jt)\,R_{ZZ}(2Jt)}
\label{eq:Uij-gates}
\end{equation}
verificato numericamente (errore $\sim10^{-15}$ rispetto a \texttt{scipy.linalg.expm}, coerente coi self-test 3/5 di \texttt{\seqsplit{trotter\_trimero\_catena.py}}).

\section{Fattorizzazione esatta di $H_0$ = campo + scambio totale}
\label{sec:H0-fact}

\subsection{Il campo commuta con ciascun singolo legame, esattamente}

Con $S_z^{tot}=\sum_iZ_i$, la stessa dimostrazione per calcolo esplicito dell'anello (Eq.~\eqref{eq:Sztot-comm-bond} di quel documento) vale \emph{qubit per qubit}, senza dipendere dal numero di legami del sistema:
\begin{equation}
[S_z^{tot},\,\vec\sigma_i\cdot\vec\sigma_j] = 0 \qquad \forall\,(i,j).
\label{eq:Sztot-comm-bond}
\end{equation}
Verificato: $\|[S_z^{tot},h_{12}]\|=\|[S_z^{tot},h_{23}]\|=0$ a precisione macchina (self-test 1 di \texttt{\seqsplit{trotter\_trimero\_catena.py}}), quindi $[S_z^{tot},H_{ex}]=0$.

\subsection{Fattorizzazione esatta di $e^{-iH_0t}$}

Per l'argomento generale di Sezione~\ref{sec:BCH}:
\begin{equation}
\boxed{e^{-iH_0t} = e^{-ibt\,S_z^{tot}}\; e^{-iH_{ex}t} = R_z^{(1)}(2bt)\,R_z^{(2)}(2bt)\,R_z^{(3)}(2bt)\;e^{-iH_{ex}t}}
\label{eq:H0-factor}
\end{equation}
\textbf{esatto, senza alcuna approssimazione}, indipendentemente da $J,b$.

\begin{center}
\fbox{\parbox{0.9\linewidth}{\textbf{Un limite netto, da non oltrepassare.} L'Eq.~\eqref{eq:H0-factor} riguarda \emph{esclusivamente} la coppia (campo, scambio). Non implica, e non va usata per suggerire, alcuna proprietà di $[H_{ex},H_{DM}]$: sono condizioni algebricamente indipendenti (Sezione~\ref{sec:non-sequitur}).}}
\end{center}

\section{I due legami di scambio non commutano: l'operatore di chiralità}
\label{sec:chirality}

\subsection{Perché serve un Trotter interno, anche con due soli legami}

$H_{ex}=J\,\vec\sigma_1\!\cdot\!\vec\sigma_2+J\,\vec\sigma_2\!\cdot\!\vec\sigma_3$ è somma di due operatori a due qubit che condividono il qubit 2. Come nell'anello, i due termini non commutano in generale: qui però c'è un solo commutatore da calcolare, non tre.

\subsection{Derivazione del commutatore}

Identico calcolo dell'anello (Eq.~\eqref{eq:chirality-comm} di quel documento), applicato all'unica coppia di legami condivisi presente nella catena:
\begin{equation}
\boxed{[\vec\sigma_1\!\cdot\!\vec\sigma_2,\,\vec\sigma_2\!\cdot\!\vec\sigma_3] = -2i\,\chi, \qquad \chi\equiv\vec\sigma_1\cdot(\vec\sigma_2\times\vec\sigma_3)}
\label{eq:chirality-comm}
\end{equation}
\textbf{Verificato numericamente}: $\|[h_{12},h_{23}]+2i\chi\|=0$ a precisione macchina (self-test 3 di \texttt{\seqsplit{trotter\_trimero\_catena.py}}), zero esatto, non un residuo di arrotondamento.

$\chi$ è lo stesso operatore di chiralità di spin scalare identificato per l'anello (Wen--Wilczek--Zee 1989). La Sezione~\ref{sec:frustrazione-nota} discute perché questo operatore compare qui \textbf{nonostante la catena non abbia alcun ciclo chiuso e nessuna frustrazione}: l'identità dell'Eq.~\eqref{eq:chirality-comm} è un fatto puramente algebrico su due operatori di Pauli che condividono un qubit, non un fatto sulla geometria del sistema.

\subsection{Conseguenza: $H_{ex}$ non fattorizza esattamente nei due legami}

Poiché $\chi\neq0$ come operatore, $e^{-iH_{ex}t}\neq U_{23}(t)U_{12}(t)$ esattamente: solo approssimativamente, con errore quantificato in Sezione~\ref{sec:trotter-error}. \textbf{Livello 1} di Trotter interno.

\subsection{Il segno dipende dall'ordine: una possibilità che l'anello non ha}
\label{sec:sign-order}

Nell'anello, la somma pesata dei tre commutatori ciclici collassa a un unico multiplo di $\chi$ (Eq.~28 del documento anello), indipendentemente dall'ordine scelto per i tre fattori del prodotto. Qui, con un solo commutatore, l'ordine dei due fattori determina direttamente il segno dell'errore:
\begin{align}
e^{-iJt h_{23}}\,e^{-iJt h_{12}} &= \exp\!\big[-it(H_{ex}-tJ^2\chi)+O(t^3)\big], \label{eq:order-12-23}\\
e^{-iJt h_{12}}\,e^{-iJt h_{23}} &= \exp\!\big[-it(H_{ex}+tJ^2\chi)+O(t^3)\big], \label{eq:order-23-12}
\end{align}
\textbf{verificato numericamente per proiezione diretta sul sottospazio di $\chi$}: coefficiente $s=+1.000$ per l'ordine $(1,2)\to(2,3)$ applicato nel circuito (che realizza il prodotto dell'Eq.~\eqref{eq:order-23-12}, poiché in Qiskit il primo gate scritto agisce per primo, quindi sta a destra nel prodotto di operatori), $s=-1.000$ per l'ordine invertito. Questa dipendenza dal segno è la base della Sezione~\ref{sec:alternanza}.

\section{Decomposizione esatta di un singolo legame DM}
\label{sec:U2-bond}

Identica, qubit per qubit, alla Sezione~5 dell'anello: per una coppia $(i,j)$ con coefficiente $D_{ij}$, $D_{ij}(X_iZ_j-Z_iX_j)$, con $P_1=X_iZ_j,P_2=Z_iX_j$ che commutano ($P_1P_2=P_2P_1=Y_iY_j$), quindi
\begin{equation}
\boxed{U_{ij}^{DM}(t) = \big[(I\otimes H)_j\,R_{ZZ}(-2D_{ij}t)\,(I\otimes H)_j\big]\cdot\big[(H\otimes I)_i\,R_{ZZ}(2D_{ij}t)\,(H\otimes I)_i\big]}
\label{eq:UDM-bond}
\end{equation}
esatto, verificato numericamente (self-test 5).

\section{I due legami DM non commutano: l'operatore $\Theta$}
\label{sec:DM-noncomm}

\subsection{Derivazione del commutatore}

Con $D_{12}=+D$, $D_{23}=-D$ (Sezione~\ref{sec:hamiltoniana}), e ponendo $d_{ij}\equiv X_iZ_j-Z_iX_j$ (coefficiente unitario), lo stesso calcolo dell'anello (Eq.~\eqref{eq:DM-comm} di quel documento, con l'indice 3 al posto di 1 poiché qui manca il legame $(3,1)$) dà:
\begin{equation}
\boxed{[d_{12},\,-d_{23}] = -2i\,\Theta, \qquad \Theta\equiv X_1Y_2Z_3-Z_1Y_2X_3}
\label{eq:DM-comm}
\end{equation}
\textbf{Verificato numericamente}: $\|[d_{12},-d_{23}]+2i\Theta\|=0$ esatto (decomposizione Pauli diretta del commutatore, due sole label non nulle). $\Theta$ è hermitiano (verificato) e \textbf{dispari sotto $P_{13}$}: la riflessione scambia i due addendi di $H_{DM}$, quindi ribalta il segno del loro commutatore -- coerente con la struttura di simmetria della Sezione~\ref{sec:hamiltoniana}.

\subsection{$\Theta$ non è proporzionale a $\chi$}

A differenza del caso a tre legami dell'anello (dove la domanda è se i tre commutatori DM collassino a un unico operatore, risposta: no, restano tre distinti, Sezione~7 di quel documento), qui la domanda più diretta è se l'unico commutatore DM ($\Theta$) sia proporzionale all'unico commutatore di scambio ($\chi$). \textbf{Verificato numericamente: no.} Livello 1 e livello 2 generano strutture d'errore indipendenti; il livello 2 non è assorbibile in una ridefinizione del livello 1.

\subsection{Conseguenza: anche $H_{DM}$ non fattorizza esattamente}

$e^{-iH_{DM}t}\neq U_{23}^{DM}(t)U_{12}^{DM}(t)$ esattamente per $D\neq0$: \textbf{livello 2} di Trotter interno, stessa struttura qualitativa dell'anello (un secondo livello non previsto dal compito originale sul dimero), qui con un solo commutatore anziché tre.

\section{Perché $[A,B]=0$ implica fattorizzazione esatta}
\label{sec:BCH}

Identico, parola per parola, alla Sezione~8 di \texttt{\seqsplit{quantum\_simulation\_trimero\_anello\_teoria.tex}} e a \texttt{\seqsplit{quantum\_simulation\_dimero\_teoria.tex}} \S5: per BCH, $[A,B]=0\Rightarrow e^Ae^B=e^{A+B}$ esattamente (ogni commutatore annidato si annulla per induzione). Si applica a: i blocchi di singolo legame (Sezioni~\ref{sec:bond-exact},~\ref{sec:U2-bond}), la fattorizzazione di $H_0$ (Sezione~\ref{sec:H0-fact}). \textbf{Non} si applica a: la somma dei due legami di scambio (Sezione~\ref{sec:chirality}), la somma dei due legami DM (Sezione~\ref{sec:DM-noncomm}), né alla coppia $H_0,H_{DM}$ -- e su quest'ultimo punto la Sezione seguente corregge un errore fatto in precedenza per l'anello.

\section{Il livello 0 esterno: cosa NON implica la fattorizzazione di $H_0$}
\label{sec:non-sequitur}

\subsection{L'errore da non ripetere}

La documentazione del modulo Trotter dell'anello giustificava l'errore Trotter esterno affermando che $H_{ex}$ non vi contribuisce «perché commuta esattamente col campo» ($[S_z^{tot},H_{ex}]=0$, Eq.~\eqref{eq:Sztot-comm-bond}). \textbf{È un non-sequitur}: $[S_z^{tot},H_{ex}]=0$ e $[H_{ex},H_{DM}]=0$ sono condizioni logicamente indipendenti; la prima non implica in alcun modo la seconda.

\subsection{Verifica numerica diretta}

A $J=1,b=0.7,D=0.3$:
\begin{equation}
\|[H_{ex},H_{DM}]\| = 10.182338, \qquad \|[S_z^{tot},H_{DM}]\| = 3.394113,
\end{equation}
norma di Frobenius. \textbf{$H_{ex}$ è il contributo dominante}, non trascurabile: circa il triplo del campo. Il livello 0 esterno non si scompone in un termine nullo (scambio) più uno residuo (campo): entrambi contribuiscono, e lo scambio di più.

\subsection{L'errore corretto, per disuguaglianza triangolare diretta}

Senza assumere alcuna cancellazione,
\begin{equation}
\big\|[H_0,H_{DM}]\big\| = \big\|b[S_z^{tot},H_{DM}]+[H_{ex},H_{DM}]\big\| \le b\,\|[S_z^{tot},H_{DM}]\|+\|[H_{ex},H_{DM}]\|.
\label{eq:level0-bound}
\end{equation}
Verificato numericamente allo stesso punto: $\|[H_0,H_{DM}]\|=10.4559$, contro il limite superiore $b\cdot3.3941+10.1823=12.558$ dell'Eq.~\eqref{eq:level0-bound} -- coerente, il limite non è saturato (i due termini non sono allineati in fase), ma resta un limite valido.

\section{Struttura a tre livelli dell'errore di Trotter}
\label{sec:trotter-error}

Il circuito (\texttt{\seqsplit{trotter\_trimero\_catena.py}}) approssima un passo come
\begin{align}
U_{\rm step}(\tau) &= V_{DM}(\tau)\;e^{-ib\tau S_z^{tot}}\;V_{ex}(\tau), \label{eq:Ustep-def}\\
V_{ex}(\tau)&=U_{23}(\tau)\,U_{12}(\tau), \qquad V_{DM}(\tau)=U_{23}^{DM}(\tau)\,U_{12}^{DM}(\tau), \notag
\end{align}
mentre l'evoluzione esatta è $U_{\rm esatto}(\tau)=e^{-i(H_0+H_{DM})\tau}$.

\subsection{Livello 1: errore di legame per $H_{ex}$}

Per due operatori, espandendo in serie fino a $O(\tau^2)$:
\begin{equation}
e^{-i(A+B)\tau}-e^{-iB\tau}e^{-iA\tau} = -\frac{\tau^2}2[A,B]+O(\tau^3).
\label{eq:2term-BCH}
\end{equation}
Con $A=Jh_{12}$, $B=Jh_{23}$ (ordine del circuito: $h_{12}$ applicato per primo, quindi a destra), usando l'Eq.~\eqref{eq:chirality-comm} con il segno di Sezione~\ref{sec:sign-order}:
\begin{equation}
\boxed{e^{-iH_{ex}\tau}-V_{ex}(\tau) = i\,\tau^2\,J^2\,\chi + O(\tau^3)}
\label{eq:error-Hex-final}
\end{equation}
A differenza dell'anello, qui \textbf{non c'è alcuna cancellazione di termini incrociati da invocare}: con due soli legami c'è un solo commutatore, non tre, quindi nulla da cancellare. La struttura è più semplice, non un caso particolare della formula dell'anello.

\subsection{Livello 2: errore di legame per $H_{DM}$}

Stesso metodo, con $A=Dd_{12}$, $B=-Dd_{23}$, usando l'Eq.~\eqref{eq:DM-comm}:
\begin{equation}
\boxed{e^{-iH_{DM}\tau}-V_{DM}(\tau) = i\,\tau^2\,D^2\,\Theta + O(\tau^3)}
\label{eq:error-HDM-final}
\end{equation}

\subsection{Livello 0 esterno: errore fra $H_0$ e $H_{DM}$}

Stesso argomento a due termini di Eq.~\eqref{eq:2term-BCH}:
\begin{equation}
e^{-i(H_0+H_{DM})\tau}-e^{-iH_{DM}\tau}e^{-iH_0\tau} = -\frac{\tau^2}2[H_0,H_{DM}]+O(\tau^3),
\label{eq:error-outer}
\end{equation}
con $[H_0,H_{DM}]\neq0$ quantificato in Eq.~\eqref{eq:level0-bound} -- correttamente, non per un presunto annullamento del contributo di $H_{ex}$.

\subsection{Composizione dei tre livelli: perché l'errore totale resta $O(\tau^2)$}

Identica, per struttura, alla Sezione~9.5 dell'anello: definendo $\tilde U_0(\tau)=e^{-ib\tau S_z^{tot}}V_{ex}(\tau)$,
\begin{equation}
\big\|U_{\rm esatto}(\tau)-U_{\rm step}(\tau)\big\| \le \|\varepsilon_0(\tau)\|+\|\varepsilon_1(\tau)\|+\|\varepsilon_2(\tau)\|,
\label{eq:triangle-step}
\end{equation}
con i tre termini a destra rispettivamente $O(\tau^2)$ dalle Eq.~\eqref{eq:error-outer},~\eqref{eq:error-Hex-final},~\eqref{eq:error-HDM-final}, quindi
\begin{equation}
\boxed{\big\|U_{\rm esatto}(\tau)-U_{\rm step}(\tau)\big\| = O(\tau^2)}.
\label{eq:total-error-bound}
\end{equation}
\textbf{Nota di cautela, identica all'anello}: la disuguaglianza triangolare limita l'errore dall'alto, non ne fissa il segno relativo; i tre contributi possono parzialmente cancellarsi o sommarsi (Sezione~\ref{sec:alternanza} sfrutta esplicitamente questo fatto per il livello 1).

\subsection{Errore totale su $N$ passi}

Identico all'anello e al dimero: $\|U_{\rm esatto}(t)-U_{\rm step}(\tau)^N\|\lesssim N\cdot O(\tau^2)=O(t^2/N)$, $\tau=t/N$.

\section{Alternanza dell'ordine dei bond: un fenomeno specifico della catena}
\label{sec:alternanza}

\subsection{L'idea}

Poiché il segno del termine $\chi$ nell'Eq.~\eqref{eq:error-Hex-final} dipende dall'ordine di applicazione dei due bond (Sezione~\ref{sec:sign-order}), alternare $(1,2)\!\to\!(2,3)$ e $(2,3)\!\to\!(1,2)$ su passi consecutivi cancella il termine di livello 1 al leading order, \textbf{a parità esatta di gate}: nessun costo aggiuntivo, un semplice riordino.

\subsection{Perché l'anello non ha un analogo altrettanto pulito}

Con tre legami che condividono i qubit ciclicamente, un'alternanza dell'ordine nell'anello non isola un singolo segno da invertire nello stesso modo: la somma dei tre commutatori ciclici già collassa a un solo termine ($-2iJ'^2\chi$, indipendente dall'ordine per la ciclicità della somma), quindi non c'è un secondo ordinamento "opposto" da alternare con lo stesso effetto di cancellazione pulita. La possibilità della Sezione presente è specifica della topologia a due soli legami.

\subsection{Quantificazione e suo limite}

Per due passi consecutivi con ordine invertito, l'errore di livello 1 di un doppio passo si cancella al primo ordine in $\tau$, lasciando un residuo $O(\tau^3)$ per coppia di passi anziché $O(\tau^2)$ per passo singolo: guadagno $O(1/N^2)\to O(1/N^4)$ per il solo livello 1, \textbf{verificato numericamente} (stato casuale, $D=0$: rapporto dell'infedeltà $\times16.0$ raddoppiando $N$, contro $\times4.0$ dell'ordine fisso).

\textbf{Con $D\neq0$ il guadagno è reale ma parziale}: l'alternanza cancella solo il termine proporzionale a $\chi$ (livello 1); il livello 2 (Eq.~\eqref{eq:error-HDM-final}) e i termini incrociati del livello 0 esterno (Eq.~\eqref{eq:level0-bound}) restano $O(\tau^2)$ per passo e tornano a dominare per $N$ grande. Verificato: a $D=0.3$ il rapporto migliora fino a $\times9.9$ ma degrada verso $\times5.1$ all'aumentare di $N$, non $\times16$ stabile.

\section{Riepilogo dei risultati}

\begin{table}[h]
\centering
\footnotesize
\begin{tabular}{@{}p{3.0cm}p{7.6cm}p{3.4cm}@{}}
\toprule
\textbf{Blocco} & \textbf{Espressione} & \textbf{Natura} \\
\midrule
$U_{ij}(t)$ & $R_{XX}(2Jt)R_{YY}(2Jt)R_{ZZ}(2Jt)$ & Esatta \\[4pt]
$U_{ij}^{DM}(t)$ & sandwich di Hadamard, Eq.~\ref{eq:UDM-bond} & Esatta \\[4pt]
$e^{-iH_0t}$ & $R_z^{(1)}(2bt)R_z^{(2)}(2bt)R_z^{(3)}(2bt)\cdot e^{-iH_{ex}t}$ & Esatta \\[4pt]
$e^{-iH_{ex}t}$ (2 legami) & $U_{23}(t)U_{12}(t)$ & Approssimata, $i\tau^2J^2\chi+O(\tau^3)$ \\[6pt]
$e^{-iH_{DM}t}$ (2 legami) & $U_{23}^{DM}(t)U_{12}^{DM}(t)$ & Approssimata, $i\tau^2D^2\Theta+O(\tau^3)$ \\[6pt]
$e^{-i(H_0+H_{DM})t}$ & $\big(V_{DM}(\tau)e^{-ib\tau S_z^{tot}}V_{ex}(\tau)\big)^N$ & Approssimata, $O(t^2/N)$ \\
\bottomrule
\end{tabular}
\caption{Riepilogo dei blocchi di evoluzione temporale per il trimero a catena aperta.}
\end{table}

\section*{Riferimenti bibliografici}
\addcontentsline{toc}{section}{Riferimenti bibliografici}

\begin{itemize}
\item H.\,F.\ Trotter, \emph{On the product of semi-groups of operators}, Proc.\ Am.\ Math.\ Soc.\ \textbf{10}, 545--551 (1959).
\item M.\ Suzuki, \emph{Generalized Trotter's formula...}, Commun.\ Math.\ Phys.\ \textbf{51}, 183--190 (1976).
\item S.\ Lloyd, \emph{Universal quantum simulators}, Science \textbf{273}, 1073--1078 (1996).
\item X.\,G.\ Wen, F.\ Wilczek, A.\ Zee, \emph{Chiral spin states and superconductivity}, Phys.\ Rev.\ B \textbf{39}, 11413 (1989).
\item M.\ Crippa et al., \emph{Simulating Static and Dynamic Properties of Magnetic Molecules with Prototype Quantum Computers}, Magnetochemistry \textbf{7}, 117 (2021).
\end{itemize}

\end{document}
